% !PRINTED 647
\weberpage{647}
From this it follows:

2. \gesperrt{If a product $\alpha\beta$ is divisible by $\pi$, then at least one of the two factors, $\alpha$ or $\beta$, is divisible by $\pi$.}

For if neither $\alpha$ nor $\beta$ is divisible by $\pi$, then both are relatively prime to $\pi$, and by \S\,181, 3. $\pi$ is also relatively prime to $\alpha\beta$.

3. \gesperrt{Every integral functional $\overline{\omega}$ different from zero is divisible by a prime functional.}

For if $\overline{\omega}$ is not itself a prime functional, then it is divisible by a functional $\alpha$ different from $\omega$. % [sic] printed without the expected bar over omega
If
\[
\overline{\omega} = \alpha \overline{\omega}',
\]
then $\overline{\omega}'$ is not a unit, and
\[
N_a(\overline{\omega}) = N_a(\alpha)\, N_a(\overline{\omega}').
\]

The absolute norms occurring here are integral functions in $z$, and the degree of $N_a(\overline{\omega}')$ is smaller than the degree of $N_a(\omega)$. % [sic] printed without the expected bar over omega
If $\overline{\omega}'$ is still not a prime functional, one can continue in this way and must finally arrive at a prime factor $\pi$ of $\overline{\omega}$.

In this way one argues further as in Bd.\,II, \S\,158, except that here the \gesperrt{height of the degree} of integral rational functions of the variable $z$ takes the place of the \gesperrt{quantity} of the integral rational numbers used there. Thus one also obtains the theorems:

4. \gesperrt{Every integral functional $\overline{\omega}$ in $\overline{\Omega}$ that is not a unit can be decomposed into a finite number of prime factors, and indeed in only one way, if associated functionals are regarded as not different.}

In what follows we must make use of the theorem that an integral rational function of $z$ with numerical coefficients is decomposable into linear factors; in other words, we must presuppose the fundamental theorem of algebra on the existence of roots. From this there follows first of all:

5. \gesperrt{The integral rational function of $z$ of lowest degree that is divisible by a prime functional $\pi$ is a linear function $z - c$.}

By \S\,180, 9. there do exist integral % [sic] printed "ganz rationale" without final -e, elsewhere Weber writes "ganze rationale"
rational functions of $z$ that are divisible by $\pi$. If one decomposes such a func-
%CONT>>
